% Use the official AAAI style when aaai2027.sty is available.
% Otherwise, fall back to article mode for drafting.
\IfFileExists{aaai2027.sty}{
    \documentclass[letterpaper]{article}
    \usepackage[submission]{aaai2027}
    \newif\ifaaaistyle
    \aaaistyletrue
}{
    \documentclass[10pt,letterpaper]{article}
    \newif\ifaaaistyle
    \aaaistylefalse
}

% ============================================================
% Fonts and basic formatting
% ============================================================
\ifaaaistyle
    \usepackage{times}
    \usepackage{helvet}
    \usepackage{courier}
\else
    \usepackage[margin=1in]{geometry}
    \usepackage{mathpazo}
\fi

% ============================================================
% Mathematics
% ============================================================
\usepackage{amsmath}
\usepackage{amssymb}
\usepackage{amsthm}
\usepackage{mathtools}
\usepackage{dsfont}
\usepackage{mathrsfs}
\usepackage{dsfont}

% ============================================================
% Figures and tables
% ============================================================
\usepackage{graphicx}
\usepackage{booktabs}
\usepackage{array}
\usepackage{siunitx}

% ============================================================
% Lists, algorithms, and code
% ============================================================
\usepackage{enumitem}
\usepackage{algorithm}
\usepackage{algorithmic}
\usepackage{listings}

% ============================================================
% References and URLs
% ============================================================
\usepackage{url}
\usepackage{natbib}

% Do not load hyperref for the AAAI submission unless the
% official author kit explicitly permits it.
\ifaaaistyle\else
    \usepackage[hidelinks]{hyperref}
\fi

% AAAI papers normally use unnumbered section headings.
\setcounter{secnumdepth}{0}

% ============================================================
% General mathematical commands
% ============================================================
\newcommand{\R}{\mathbb{R}}
\newcommand{\N}{\mathbb{N}}
\newcommand{\Z}{\mathbb{Z}}
\newcommand{\Q}{\mathbb{Q}}
\newcommand{\E}{\mathbb{E}}
\newcommand{\Prob}{\mathbb{P}}
\newcommand{\1}{\mathds{1}}
\newcommand{\xeta}{\xi, \eta}

\newcommand{\calF}{\mathcal{F}}
\newcommand{\bigO}{\mathcal{O}}

\newcommand{\kpar}{\kappa}
\newcommand{\cs}{\sigma}
\newcommand{\mo}{\mu}
\newcommand{\lm}{\text{log\,MAE}}

\DeclareMathOperator{\Var}{Var}
\DeclareMathOperator{\Cov}{Cov}
\DeclareMathOperator{\Aut}{Aut}
\DeclareMathOperator{\Gal}{Gal}
\DeclareMathOperator{\lcmop}{lcm}
\DeclareMathOperator{\argmin}{arg\,min}
\DeclareMathOperator{\argmax}{arg\,max}

% ============================================================
% Paper-specific commands
% ============================================================
\newcommand{\parentIntensity}{\kappa}
\newcommand{\clusterScale}{\sigma}
\newcommand{\meanOffspring}{\mu}
\newcommand{\logMAE}{\operatorname{logMAE}}

% ============================================================
% Draft annotations
% ============================================================
\newif\ifdraft
\drafttrue

\ifdraft
    \usepackage{xcolor}
    \newcommand{\todo}[1]{\textcolor{red}{[TODO: #1]}}
\else
    \newcommand{\todo}[1]{}
\fi

% ============================================================
% Theorem environments
% ============================================================
\theoremstyle{plain}
\newtheorem{theorem}{Theorem}
\newtheorem{proposition}[theorem]{Proposition}
\newtheorem{lemma}[theorem]{Lemma}
\newtheorem{claim}[theorem]{Claim}
\newtheorem{corollary}[theorem]{Corollary}

\theoremstyle{definition}
\newtheorem{definition}[theorem]{Definition}
\newtheorem{assumption}[theorem]{Assumption}
\newtheorem{problem}[theorem]{Problem}

\theoremstyle{remark}
\newtheorem{remark}[theorem]{Remark}

% ============================================================
% Title and author information
% ============================================================
\title{Deep Learning for Parameter Estimation in Homogeneous Thomas Processes}

\ifaaaistyle
    % Anonymous submission version
    \author{Anonymous Submission}

    % Replace the anonymous author block with the following for
    % the camera-ready version:
    %
    % \author{
    %     Flavio Gualtieri\textsuperscript{\rm 1}
    % }
    %
    % \affiliations{
    %     \textsuperscript{\rm 1}Institution Name\\
    %     City, Country\\
    %     email@example.com
    % }
\else
    \author{Flavio Gualtieri}
    \date{\today}
\fi

\begin{document}

\maketitle

\begin{abstract}
We propose a novel deep learning estimator for the three parameters of a homogeneous Thomas point process: parent intensity $\kpar$, cluster scale $\cs$, and mean offspring $\mo$. The estimator is comprised of an encoder that computes a $128$-dimensional feature vector, which is in turn fed to a FCNN with three output nodes corresponding to the three parameters. We present a comparison of various model features - raw point cloud, parwise distance CDF, $H_0$/$H_1$ persistence image, $H_0$/$H_1$ Betti curve - showing that features constructed from the persistent homology of the point cloud perform significantly better than more classical and second-order features.
\end{abstract}


\section{Introduction}
The study of spatial point processes (SPPs) has a variety of important applications spanning astronomy, cell biology, and epidemiology. Many natural phenomena can be modelled using point processes - such as galaxy clustering - so one can hope to leverage an understanding of SPPs to obtain new insights.

The ability to infer the parameters that describe a given point process is of particular interest becasue a wealth of summary statistics and properties become available once these parameters are known.

A number of parameter estimation methods have been proposed, typically relying on closed-form expressions for summary staistics such as the $K$ and $g$ functions. Minimum contrast on $K$ or $g$ (Diggle 2003). Palm likelihood (Tanaka, Ogata \& Stoyan 2008). Bayesian and MCMC.

These methods have limitations we hope to address. Reliance on hand-picked summary statistic (and by extension a closed-form expression for this statistic), hyperparameter tuning sensitivity.


\section{Background}
We provide a basic introduction to spatial point processes and notions from persistent homology used in this paper.

\subsection{The Thomas Point Process}
A $d$-dimensional spatial point process $X$ on a subset $W \subseteq \R^d$ is equipped with an intensity function $\rho(u): \R^d \rightarrow \R^+$ such that

    \[
    \E[N(W)] = \int_W \rho(u) du,
    \]
    
where $N(W) = |W \cap X|$ denotes the number of points of $X$ in $W$. When $\rho(u) = \lambda \in \R$ we say that $X$ is \emph{homogeneous}.

    \begin{definition}[Neimann-Scott process]
        For a Neimann-Scott process with Poisson-distributed offspring we have:
        
            \[
            \rho_{NS}(u) = \mo \sum_{p \in P} k(u - p),
            \]
            
        where $P$ is a Poisson-distributed \emph{parent} point process with $\rho_P(u) = \kpar$, and $k$ is a kernel that describes how the \emph{offspring} point are spatially distributed around their parents. We therefore identify the parameters of parent and offspring intensities, as well as any additional parameters introduced by the kernel $k$.
    \end{definition}

A homogeneous Thomas process is a special type of Neimann-Scott process where $k$ is an isotropic Gaussian kernel $\mathcal{N}(0, \cs^2I)$, and therefore has parameters $\theta = (\kpar, \mo, \cs)$. Computationally speaking, we may construct a Thomas point process on a $d$-dimensional unit box $W = [0, 1]^d$ as follows:

    \begin{enumerate}
        \item Expand $W$ by an edge buffer $b = 4\cs$ to avoid boundary truncation of clusters, to obtain $W_{exp}$.
        \item Sample a homogeneous Poisson parent process $P = \{p_i\}_{i=1}^{N_p}$ with intensity $\rho_P(u) = \kpar$ on $W_{exp}$, so $N_p \sim Po(\kpar \cdot |W_{exp}|)$.
        \item Each $p_i \in P$ produces $K_p \sim Po(\mo)$ offspring.
        \item Displace offspring from their respective parents with i.i.d. isotropic Gaussian noise, $x_{ij} = p_i + \epsilon_{ij}$, where $\epsilon_{ij} \sim \mathcal{N}(0, \sigma^2I_d)$.
        \item Truncate $W_{exp}$ back to $W$, discarding points in $\R^d \setminus W$ and retaining only the remaining offspring points to obtain $X_{Thomas} = \{x_k\}_{k=1}^{N_T} \subset [0, 1]^d$.
    \end{enumerate}


\subsection{Summary Statistics}
A wealth of summary statistics exist to describe SPPs. In this paper we focus on the $g$ and $K$ functions because they are the most widely used in parameter estimation, but these are no means the only ones. We begin by defining the second-order product density $\rho^{(2)}(u, v)$.

    \begin{definition}[Second-order product density]
        Let $X$ denote a point process on $\R^d$ with $\rho(u)$ and $\xeta \in X$. The function $\rho^{(2)}(u, v):  \R^d \times \R^d \rightarrow \R^+$ is the second-order product density if:

            \[
            \alpha^{(2)}(C) = \int \int \1[(\xeta) \in C]\rho^{(2)}(\xeta) d\xi d\eta, \quad C \subseteq \R^d \times \R^d,
            \]

        where $\alpha^{(2)}(C): R^d \times \R^d \rightarrow \R^+$ is the second-order factorial moment measure:

            \[
            \alpha^{(2)}(C) = \E\sum_{\xeta \in X}^{\neq} \1[(\xeta) \in C], \quad C \subseteq \R^d \times \R^d.
            \]
    \end{definition}

The second-order product density $\rho^{(2)}(u, v)$ describes the probability of jointly observing a pair of points from $X$ in balls $b_{\epsilon}(\xi)$, $b_{\epsilon}(\eta)$ with volumes $d\xi$, $d\eta$. The pair correlation function $g(\xeta): R^d \times R^d \rightarrow \R^+$ is defined

    \[
    g(\xeta) = \frac{\rho^{(2)}(\xeta)}{\rho(\xi)\rho(\eta)}.
    \]

It is clear that for a homogeneous Poisson process we have $g = 1$. Intuitively, $g$ provides information about whether an SPP is attractive ($g > 1$), or repulsive ($g < 1$). We now define the $K$ function.

    \begin{definition}[$K$-function]
        For a stationary and isotropic point process $X$ on $\R^d$, the $K$-function $K(r): \R \rightarrow R^+$ is defined

            \[
            K(r) = \int_{b_r(0)} g(\xi) d\xi
            \]
    \end{definition}

Intuitively, $K(r)$ provides a measure of how many further points we expect to see at a distance $r$ from the origin.

\subsection{Persistent Homology}
We now turn out attention to some basic notions from persistent homology (PH).

    \begin{definition}[Vietoris-Rips complex]
        Let $X = \{x_i\}_{i=1}^N$ where $x_i \in \R^d$ denote a point cloud. We equip $X$ with a Vietoris-Rips (VR) complex by defining:
            
            \[
            VR_t(X) = \{\sigma | b_t(x_i) \cap b_t(x_j) \neq \emptyset \forall x_i, x_j \in \sigma\}
            \]
    \end{definition}

From the notion of a VR complex the notion of VR filtration follows naturally, as one can easily see that $VR_s(X) \subseteq VR_t(X)$ whenever $s \leq t$.


- Define birth-death
- Define persistence diagram
- Define persitence image
- Define betti curve

\section{Model}
We explain the model and it's features

\subsection{Overall Architecture}
- Define FNCC
- Define encoders
- Explain specifics of this model's head (with three output nodes)

\subsection{Features}
Explain each feature and it's encdoer

\subsection{Raw Point Cloud}
- Point-net style encoder
- Explain importance of persumation invariance

\subsection{Pairwise Distance CDF}

\subsection{Persistence Image}
- Explain how it is computed from diagrams
- Explain calibration procedure
- Explain CNN encoder

\subsection{Betti Curve}
- Sample 128 points on betti curve


\section{Results}
Explain results





\end{document}